% Options for packages loaded elsewhere
\PassOptionsToPackage{unicode}{hyperref}
\PassOptionsToPackage{hyphens}{url}
\documentclass[
]{article}
\usepackage{xcolor}
\usepackage{amsmath,amssymb}
\setcounter{secnumdepth}{-\maxdimen} % remove section numbering
\usepackage{iftex}
\ifPDFTeX
  \usepackage[T1]{fontenc}
  \usepackage[utf8]{inputenc}
  \usepackage{textcomp} % provide euro and other symbols
\else % if luatex or xetex
  \usepackage{unicode-math} % this also loads fontspec
  \defaultfontfeatures{Scale=MatchLowercase}
  \defaultfontfeatures[\rmfamily]{Ligatures=TeX,Scale=1}
\fi
\usepackage{lmodern}
\ifPDFTeX\else
  % xetex/luatex font selection
\fi
% Use upquote if available, for straight quotes in verbatim environments
\IfFileExists{upquote.sty}{\usepackage{upquote}}{}
\IfFileExists{microtype.sty}{% use microtype if available
  \usepackage[]{microtype}
  \UseMicrotypeSet[protrusion]{basicmath} % disable protrusion for tt fonts
}{}
\makeatletter
\@ifundefined{KOMAClassName}{% if non-KOMA class
  \IfFileExists{parskip.sty}{%
    \usepackage{parskip}
  }{% else
    \setlength{\parindent}{0pt}
    \setlength{\parskip}{6pt plus 2pt minus 1pt}}
}{% if KOMA class
  \KOMAoptions{parskip=half}}
\makeatother
\usepackage{longtable,booktabs,array}
\usepackage{caption}
\captionsetup[table]{skip=6pt}
\newcounter{none} % for unnumbered tables
\usepackage{calc} % for calculating minipage widths
% Correct order of tables after \paragraph or \subparagraph
\usepackage{etoolbox}
\makeatletter
\patchcmd\longtable{\par}{\if@noskipsec\mbox{}\fi\par}{}{}
\makeatother
% Allow footnotes in longtable head/foot
\IfFileExists{footnotehyper.sty}{\usepackage{footnotehyper}}{\usepackage{footnote}}
\makesavenoteenv{longtable}
\usepackage{graphicx}
\makeatletter
\newsavebox\pandoc@box
\newcommand*\pandocbounded[1]{% scales image to fit in text height/width
  \sbox\pandoc@box{#1}%
  \Gscale@div\@tempa{\textheight}{\dimexpr\ht\pandoc@box+\dp\pandoc@box\relax}%
  \Gscale@div\@tempb{\linewidth}{\wd\pandoc@box}%
  \ifdim\@tempb\p@<\@tempa\p@\let\@tempa\@tempb\fi% select the smaller of both
  \ifdim\@tempa\p@<\p@\scalebox{\@tempa}{\usebox\pandoc@box}%
  \else\usebox{\pandoc@box}%
  \fi%
}
% Set default figure placement to htbp
\def\fps@figure{htbp}
\makeatother
\setlength{\emergencystretch}{3em} % prevent overfull lines
\providecommand{\tightlist}{%
  \setlength{\itemsep}{0pt}\setlength{\parskip}{0pt}}
\usepackage{bookmark}
\IfFileExists{xurl.sty}{\usepackage{xurl}}{} % add URL line breaks if available
\urlstyle{same}
% fallback for those not using the hyperref driver hyperxmp:
\makeatletter
\@ifundefined{xmpquote}{\newcommand{\xmpquote}[1]{#1}}{}
\makeatother
\hypersetup{
  hidelinks,
  pdfcreator={LaTeX via pandoc}}

\author{}
\date{}

\begin{document}

\section{Bab 9: Demografi \& Struktur Sosial: Guncangan Sosial dan
Pergeseran Ekonomi
Agraris}\label{bab-9-demografi-struktur-sosial-guncangan-sosial-dan-pergeseran-ekonomi-agraris}

\textbf{CELIOS --- Center of Economic and Law Studies}

\emph{Membaca tekanan demografi, intensifikasi ruang, dan transisi
ekonomi di lingkar industri ekstraktif Sulawesi.}

\begin{center}\rule{0.5\linewidth}{0.5pt}\end{center}

\subsection{Metodologi Pendekatan}\label{metodologi-pendekatan}

\textbf{Alur Kausalitas:} \texttt{Ekspansi\ Nikel} →
\texttt{Tekanan\ Demografi\ \&\ Kepadatan} →
\texttt{Pergeseran\ Struktur\ Ekonomi} →
\texttt{Beban\ Sosial-Kesehatan}

\begin{itemize}
\tightlist
\item
  \textbf{Variabel Tekanan (X):} Kabupaten prioritas industri
  ekstraktif, IUP kumulatif, porsi PDRB pertambangan dan industri
  pengolahan, serta nilai investasi PMDN provinsi.
\item
  \textbf{Variabel Dampak (Y):} Jumlah penduduk kabupaten, kepadatan
  penduduk, laju pertumbuhan penduduk, kasus DBD sebagai proxy tekanan
  kesehatan, dan pergeseran proporsi PDRB pertanian vs tambang+industri.
\end{itemize}

\textbf{Catatan Batasan:} Halaman ini tidak mengklaim data migrasi risen
tahunan langsung. Analisis migrasi dibaca sebagai \textbf{proxy tekanan
demografi} dari data populasi dan kepadatan kabupaten. Analisis
perubahan pekerjaan dibaca sebagai \textbf{pergeseran struktur ekonomi}
berbasis PDRB sektoral, bukan perpindahan individu pekerja secara
literal.

\begin{center}\rule{0.5\linewidth}{0.5pt}\end{center}

\subsection{Ketika Hilirisasi Mengubah Struktur
Masyarakat}\label{ketika-hilirisasi-mengubah-struktur-masyarakat}

Ekspansi nikel di Sulawesi bukan hanya perubahan industri, melainkan
rekayasa ulang ruang hidup. Data demografi dan ekonomi sektoral
menunjukkan bahwa kawasan yang menjadi pusat industri ekstraktif
mengalami tekanan ganda: populasi dan kepadatan meningkat, sementara
struktur ekonomi regional bergerak meninggalkan basis agraris menuju
dominasi tambang dan industri pengolahan. Di Sulawesi Tengah, provinsi
yang menjadi episentrum Morowali dan Morowali Utara, porsi PDRB sektor
pertanian turun dari \textbf{34.39\%} pada \textbf{2014} menjadi
\textbf{15.80\%} pada \textbf{2024}. Pada periode yang sama, gabungan
sektor pertambangan dan industri pengolahan melonjak dari
\textbf{15.45\%} menjadi \textbf{55.82\%}.

Perubahan ini tidak netral. Indeks pergeseran agraris-ke-industri di
Sulawesi Tengah naik dari \textbf{0.449} menjadi \textbf{3.533}, atau
sekitar \textbf{7.9 kali}. Pada level kabupaten, Morowali memperlihatkan
sinyal tekanan demografi yang tajam: pada 2020, data SIMDASI mencatat
penduduk sebesar \textbf{161.7 ribu jiwa} dengan laju pertumbuhan sumber
\textbf{4.54\%}. Angka-angka ini tidak cukup untuk menyebut migrasi
langsung secara definitif, tetapi cukup kuat sebagai proxy bahwa kawasan
industri ekstraktif mengalami tarikan penduduk dan intensifikasi ruang
yang tidak dialami merata oleh wilayah non-industri. Dengan demikian,
hilirisasi tidak hanya memindahkan bijih menjadi logam; ia juga
memindahkan beban sosial ke masyarakat lokal.

\begin{center}\rule{0.5\linewidth}{0.5pt}\end{center}

\subsection{Ringkasan Metrik Agregat}\label{ringkasan-metrik-agregat}

{\def\LTcaptype{none} % do not increment counter
\begin{longtable}[]{@{}
  >{\raggedright\arraybackslash}p{(\linewidth - 4\tabcolsep) * \real{0.3333}}
  >{\raggedright\arraybackslash}p{(\linewidth - 4\tabcolsep) * \real{0.3333}}
  >{\raggedright\arraybackslash}p{(\linewidth - 4\tabcolsep) * \real{0.3333}}@{}}
\toprule\noalign{}
\begin{minipage}[b]{\linewidth}\raggedright
Indikator
\end{minipage} & \begin{minipage}[b]{\linewidth}\raggedright
Nilai
\end{minipage} & \begin{minipage}[b]{\linewidth}\raggedright
Keterangan
\end{minipage} \\
\midrule\noalign{}
\endhead
\bottomrule\noalign{}
\endlastfoot
\textbf{Shift Index Sulteng} & \textbf{3.53×} & Rasio tambang+industri
terhadap pertanian. Makin tinggi berarti struktur ekonomi makin bergeser
meninggalkan basis agraris. Sumber: BPS SIMDASI \\
\textbf{Pertanian Sulteng Turun} & \textbf{34.4\% → 15.8\%} & Porsi PDRB
pertanian menyusut tajam, menunjukkan pelemahan basis ekonomi agraris
dalam struktur regional. Sumber: BPS SIMDASI \\
\textbf{Tambang+Industri Sulteng Naik} & \textbf{15.4\% → 55.8\%} &
Gabungan pertambangan dan industri pengolahan menjadi blok dominan dalam
ekonomi Sulawesi Tengah. Sumber: BPS SIMDASI \\
\textbf{Kabupaten Industri Ekstraktif} & \textbf{7 Kabupaten} &
Kabupaten prioritas untuk membaca tekanan demografi dan ekonomi di
lingkar industri ekstraktif Sulawesi. Sumber: Klasifikasi Fase 4 \\
\textbf{Rasio Kepadatan Industri Ekstraktif} & \textbf{0.10×} &
Perbandingan rata-rata kepadatan kabupaten industri ekstraktif terhadap
non-ekstraktif pada tahun 2024. Sumber: BPS SIMDASI \\
\textbf{Kasus DBD di Kab. Ekstraktif} & \textbf{1,378 Kasus} & Akumulasi
DBD sejak 2019 pada kabupaten prioritas industri ekstraktif sebagai
proxy tekanan kesehatan. Sumber: Profil Kesehatan/Dinkes \\
\end{longtable}
}

\begin{center}\rule{0.5\linewidth}{0.5pt}\end{center}

\subsection{9.1 Tekanan Demografi di Kabupaten Industri
Ekstraktif}\label{tekanan-demografi-di-kabupaten-industri-ekstraktif}

\textbf{Metode: Proxy Migrasi dari Time-Series Populasi Kabupaten}

Analisis ini membaca tekanan demografi melalui perubahan jumlah penduduk
kabupaten, bukan melalui data migrasi langsung. Dengan pendekatan ini,
populasi diperlakukan sebagai sinyal awal: ketika kawasan smelter tumbuh
lebih cepat dibanding pola umum wilayah sekitar, maka terdapat indikasi
tarikan penduduk, pekerja, dan aktivitas ekonomi baru yang perlu diuji
lebih lanjut. Fokus pembacaan ditempatkan pada tujuh kabupaten prioritas
smelter, yaitu \textbf{Banggai, Kolaka, Konawe, Konawe Utara, Luwu
Timur, Morowali, Morowali Utara}. Dalam window data yang tersedia,
rata-rata pertumbuhan YoY kabupaten smelter tercatat \textbf{1.95\%},
sedangkan wilayah non-smelter berada di sekitar \textbf{2.09\%}. Pada
tahun \textbf{2024}, total populasi kabupaten smelter mencapai
\textbf{1,588.2 ribu jiwa}. Angka-angka ini tidak cukup untuk menyebut
asal migran atau arah mobilitas penduduk, tetapi cukup kuat untuk
menunjukkan bahwa hilirisasi nikel menciptakan tekanan demografis yang
harus dibaca sebagai bagian dari beban sosial, bukan sekadar konsekuensi
administratif pembangunan industri.

\begin{center}\rule{0.5\linewidth}{0.5pt}\end{center}

\subsection{9.2 Intensifikasi Ruang: Kepadatan Industri Ekstraktif vs
Non-Ekstraktif}\label{intensifikasi-ruang-kepadatan-industri-ekstraktif-vs-non-ekstraktif}

\textbf{Metode: Comparative Density Analysis}

Sub-bab ini tidak mengklaim perubahan resmi desa menjadi kota karena
data klasifikasi Podes belum menjadi basis utama di halaman ini. Yang
dibaca adalah \textbf{intensifikasi ruang}, yaitu tekanan yang muncul
ketika pertumbuhan penduduk dan konsentrasi industri bertemu pada
wilayah yang sama. Rata-rata kepadatan kabupaten smelter pada
\textbf{2024} mencapai \textbf{42.7 jiwa/km²}, sedangkan kabupaten
non-smelter berada pada \textbf{438.3 jiwa/km²}. Rasio smelter terhadap
non-smelter sebesar \textbf{0.10 kali} memberi sinyal bahwa kawasan
industri membutuhkan kapasitas layanan publik yang berbeda: perumahan,
air bersih, sanitasi, transportasi, hingga fasilitas kesehatan. Dalam
kerangka D3TLH, kepadatan bukan sekadar angka demografi, melainkan
indikator apakah ruang hidup lokal sedang dipadatkan oleh proyek
ekstraktif tanpa perencanaan sosial yang sepadan. Karena itu, grafik
berikut dibaca sebagai peta awal tekanan ruang, bukan sebagai klaim
urbanisasi formal.

\begin{figure}
\centering
\pandocbounded{\includegraphics[keepaspectratio,alt={Rata-rata Kepadatan Penduduk: Kabupaten Industri Ekstraktif vs Non-Ekstraktif}]{visuals_bab9/chart_9_2_density_area.png}}
\caption{Rata-rata Kepadatan Penduduk: Kabupaten Industri Ekstraktif vs
Non-Ekstraktif}
\end{figure}

\begin{center}\rule{0.5\linewidth}{0.5pt}\end{center}

\subsection{9.3 Pergeseran Ekonomi Agraris ke Tambang dan
Industri}\label{pergeseran-ekonomi-agraris-ke-tambang-dan-industri}

\textbf{Metode: PDRB Sector Shift Index (B+C / A)}

Pergeseran pekerjaan tidak dapat diklaim hanya dari PDRB, tetapi
struktur PDRB memberi petunjuk kuat tentang arah ekonomi yang sedang
dibentuk. Di sini sektor A dibaca sebagai basis agraris, sementara
sektor B dan C dibaca sebagai blok ekstraktif-industrial: pertambangan
dan industri pengolahan. Rasio B+C terhadap A menjadi \emph{shift
index}; nilai di atas 1 berarti kontribusi tambang dan industri sudah
melampaui pertanian. Di Sulawesi Tengah, porsi pertanian turun dari
\textbf{34.39\%} menjadi \textbf{15.80\%}, sementara tambang+industri
naik dari \textbf{15.45\%} menjadi \textbf{55.82\%}. Indeksnya naik dari
\textbf{0.449} ke \textbf{3.533}, atau sekitar \textbf{7.9 kali}. Dengan
kata lain, data sektoral menunjukkan bahwa hilirisasi tidak hanya
menambah pabrik; ia mengubah pusat gravitasi ekonomi daerah, dari ruang
produksi agraris menuju rantai ekstraktif yang lebih terkonsentrasi pada
modal besar.

\begin{quote}
\textbf{Catatan Metodologi:} BPS menggabungkan Pertanian, Kehutanan, dan
Perikanan dalam Sektor A. \textbf{Perikanan Tangkap} diestimasi sebagai
\textbf{±22\% dari nilai Sektor A}, mengacu pada rata-rata proporsi
sub-sektor perikanan terhadap Sektor A di provinsi-provinsi pesisir
Sulawesi (Sumber: Statistik Perikanan BPS Sulawesi, 2016--2024). Sektor
B+C digabung menjadi satu blok ekstraktif-industrial.
\end{quote}

\paragraph{Komposisi PDRB Sektor Kunci --- Sulawesi
Tengah}\label{komposisi-pdrb-sektor-kunci-sulawesi-tengah}

\begin{figure}
\centering
\pandocbounded{\includegraphics[keepaspectratio,alt={Komposisi PDRB Sektor Kunci Sulawesi Tengah}]{visuals_bab9/chart_9_3a_sector_sulteng.png}}
\caption{Komposisi PDRB Sektor Kunci Sulawesi Tengah}
\end{figure}

\paragraph{Indeks Pergeseran Agrikultur vs Industri (B+C / A) per
Provinsi}\label{indeks-pergeseran-agrikultur-vs-industri-bc-a-per-provinsi}

\begin{figure}
\centering
\pandocbounded{\includegraphics[keepaspectratio,alt={Indeks Pergeseran Agrikultur vs Industri per Provinsi}]{visuals_bab9/chart_9_3b_shift_index.png}}
\caption{Indeks Pergeseran Agrikultur vs Industri per Provinsi}
\end{figure}

\begin{center}\rule{0.5\linewidth}{0.5pt}\end{center}

\subsection{9.4 Sintesis: Matriks Tekanan
Sosial-Ekologis}\label{sintesis-matriks-tekanan-sosial-ekologis}

\textbf{Metode: Executive Crosstab Sektor Ekonomi × Demografi ×
Kesehatan}

Matriks sintesis menggabungkan tiga lapis bukti: perubahan struktur
ekonomi, keberadaan kabupaten industri ekstraktif, dan beban DBD di
wilayah prioritas. Tujuannya bukan mengganti analisis kausal formal,
melainkan memberi ringkasan eksekutif untuk membaca provinsi mana yang
paling kuat menunjukkan kombinasi tekanan ekonomi-ekologis dan sosial.
Berdasarkan data yang sudah diproses, provinsi dengan kenaikan shift
index tertinggi adalah \textbf{Sulawesi Tengah}, dengan delta sebesar
\textbf{3.08} poin dari tahun awal ke tahun akhir. Ini berarti perubahan
struktur ekonomi tidak merata di seluruh Sulawesi; ada wilayah yang
mengalami transformasi jauh lebih tajam karena posisinya dalam rantai
nikel. Tabel berikut mengurutkan provinsi berdasarkan delta shift index,
lalu mengaitkannya dengan jumlah kabupaten industri ekstraktif dan total
DBD di wilayah prioritas. Dengan susunan ini, pembaca dapat melihat
bahwa tekanan sosial-ekologis bukan hanya soal satu indikator, melainkan
gabungan antara ekonomi yang makin ekstraktif, penduduk yang
terkonsentrasi, dan beban kesehatan yang muncul di ruang yang sama.

Uji crosstab berikut memakai unit observasi \textbf{kabupaten-tahun},
bukan provinsi-tahun. Perubahan ini penting karena tekanan sosial
terjadi di level kabupaten, sementara agregasi provinsi membuat variasi
lokal hilang dan membuat banyak tabel menjadi tidak signifikan. Variabel
X merepresentasikan intensitas ekonomi ekstraktif pada provinsi-tahun,
lalu diwariskan ke kabupaten di provinsi yang sama. Variabel Y
merepresentasikan kepadatan, populasi, pertumbuhan penduduk, kemiskinan,
dan beban DBD.

\subsubsection{Detail Uji Statistik (Chi-Square \& Odds
Ratio)}\label{detail-uji-statistik-chi-square-odds-ratio}

\textbf{Variabel Independen (X):} Shift Index Tambang+Industri /
Pertanian

\textbf{Variabel Dependen (Y):} Kepadatan Penduduk Kabupaten

\paragraph{Case Processing Summary}\label{case-processing-summary}

{\def\LTcaptype{none} % do not increment counter
\begin{longtable}[]{@{}
  >{\raggedright\arraybackslash}p{(\linewidth - 12\tabcolsep) * \real{0.1429}}
  >{\raggedright\arraybackslash}p{(\linewidth - 12\tabcolsep) * \real{0.1429}}
  >{\raggedright\arraybackslash}p{(\linewidth - 12\tabcolsep) * \real{0.1429}}
  >{\raggedright\arraybackslash}p{(\linewidth - 12\tabcolsep) * \real{0.1429}}
  >{\raggedright\arraybackslash}p{(\linewidth - 12\tabcolsep) * \real{0.1429}}
  >{\raggedright\arraybackslash}p{(\linewidth - 12\tabcolsep) * \real{0.1429}}
  >{\raggedright\arraybackslash}p{(\linewidth - 12\tabcolsep) * \real{0.1429}}@{}}
\toprule\noalign{}
\begin{minipage}[b]{\linewidth}\raggedright
\end{minipage} & \begin{minipage}[b]{\linewidth}\raggedright
Valid N
\end{minipage} & \begin{minipage}[b]{\linewidth}\raggedright
Valid \%
\end{minipage} & \begin{minipage}[b]{\linewidth}\raggedright
Missing N
\end{minipage} & \begin{minipage}[b]{\linewidth}\raggedright
Missing \%
\end{minipage} & \begin{minipage}[b]{\linewidth}\raggedright
Total N
\end{minipage} & \begin{minipage}[b]{\linewidth}\raggedright
Total \%
\end{minipage} \\
\midrule\noalign{}
\endhead
\bottomrule\noalign{}
\endlastfoot
Shift Index Tambang+Industri / Pertanian * Kepadatan Penduduk Kabupaten
& 608 & 100.0\% & 0 & 0.0\% & 608 & 100.0\% \\
\end{longtable}
}

\paragraph{Shift Index Tambang+Industri / Pertanian * Kepadatan Penduduk
Kabupaten
Crosstabulation}\label{shift-index-tambangindustri-pertanian-kepadatan-penduduk-kabupaten-crosstabulation}

{\def\LTcaptype{none} % do not increment counter
\begin{longtable}[]{@{}llll@{}}
\toprule\noalign{}
& Rendah (\textless119.5) & Tinggi (≥119.5) & Total \\
\midrule\noalign{}
\endhead
\bottomrule\noalign{}
\endlastfoot
\textbf{Rendah (\textless0.8)} Count & 123 & 181 & 304 \\
\textbf{Rendah (\textless0.8)} Expected & 152.0 & 152.0 & 304.0 \\
\textbf{Tinggi (≥0.8)} Count & 181 & 123 & 304 \\
\textbf{Tinggi (≥0.8)} Expected & 152.0 & 152.0 & 304.0 \\
\textbf{Total} Count & 304 & 304 & 608 \\
\textbf{Total} Expected & 304.0 & 304.0 & 608.0 \\
\end{longtable}
}

\paragraph{Chi-Square Tests}\label{chi-square-tests}

\textbf{Shift Index Tambang+Industri / Pertanian * Kepadatan Penduduk
Kabupaten}

{\def\LTcaptype{none} % do not increment counter
\begin{longtable}[]{@{}llll@{}}
\toprule\noalign{}
& Value & df & Asymp. Sig. (2-sided) \\
\midrule\noalign{}
\endhead
\bottomrule\noalign{}
\endlastfoot
Pearson Chi-Square & 21.375 & 1 & 0.000 \\
Likelihood Ratio & 21.502 & 1 & 0.000 \\
Linear-by-Linear Association & 22.095 & 1 & 0.000 \\
N of Valid Cases & 608 & & \\
\end{longtable}
}

\subsubsection{Ringkasan Uji Hipotesis}\label{ringkasan-uji-hipotesis}

\textbf{Result: SIGNIFIKAN}

{\def\LTcaptype{none} % do not increment counter
\begin{longtable}[]{@{}ll@{}}
\toprule\noalign{}
Parameter & Nilai \\
\midrule\noalign{}
\endhead
\bottomrule\noalign{}
\endlastfoot
P-Value & 0.0000 \\
Chi-Square & 21.375 \\
df & 1 \\
\textbf{Odds Ratio (Risk Estimate)} & \textbf{0.462} \\
\end{longtable}
}

\begin{quote}
\textbf{Interpretasi Sosial-Ekologis:}

Hasil signifikan menunjukkan bahwa intensitas ekonomi ekstraktif
memiliki asosiasi statistik dengan indikator tekanan sosial-demografis
yang dipilih. Dalam konteks ini, memperkuat pembacaan bahwa pergeseran
struktur ekonomi perlu dibaca bersama kepadatan, populasi smelter, dan
beban kesehatan.
\end{quote}

\begin{center}\rule{0.5\linewidth}{0.5pt}\end{center}

\subsubsection{Ringkasan Eksekutif Seluruh Skenario
Crosstab}\label{ringkasan-eksekutif-seluruh-skenario-crosstab}

Tabel di bawah ini merangkum hasil pengujian statistik (Chi-Square)
untuk semua kemungkinan kombinasi antara indikator Ekonomi Ekstraktif
(X) dan Tekanan Sosial-Demografis (Y) pada panel kabupaten-tahun yang
sama.

{\def\LTcaptype{none} % do not increment counter
\begin{longtable}[]{@{}
  >{\raggedright\arraybackslash}p{(\linewidth - 10\tabcolsep) * \real{0.1667}}
  >{\raggedright\arraybackslash}p{(\linewidth - 10\tabcolsep) * \real{0.1667}}
  >{\raggedright\arraybackslash}p{(\linewidth - 10\tabcolsep) * \real{0.1667}}
  >{\raggedright\arraybackslash}p{(\linewidth - 10\tabcolsep) * \real{0.1667}}
  >{\raggedright\arraybackslash}p{(\linewidth - 10\tabcolsep) * \real{0.1667}}
  >{\raggedright\arraybackslash}p{(\linewidth - 10\tabcolsep) * \real{0.1667}}@{}}
\toprule\noalign{}
\begin{minipage}[b]{\linewidth}\raggedright
Variabel Independen (X)
\end{minipage} & \begin{minipage}[b]{\linewidth}\raggedright
Variabel Dependen (Y)
\end{minipage} & \begin{minipage}[b]{\linewidth}\raggedright
Chi-Square
\end{minipage} & \begin{minipage}[b]{\linewidth}\raggedright
P-Value
\end{minipage} & \begin{minipage}[b]{\linewidth}\raggedright
Odds Ratio
\end{minipage} & \begin{minipage}[b]{\linewidth}\raggedright
Kesimpulan
\end{minipage} \\
\midrule\noalign{}
\endhead
\bottomrule\noalign{}
\endlastfoot
Shift Index Tambang+Industri / Pertanian & Kepadatan Penduduk Kabupaten
& 21.375 & 0.000 & 0.46 & 🟢 SIGNIFIKAN \\
Shift Index Tambang+Industri / Pertanian & Jumlah Penduduk Kabupaten
(ribu jiwa) & 5.533 & 0.019 & 1.49 & 🟢 SIGNIFIKAN \\
Shift Index Tambang+Industri / Pertanian & Laju Pertumbuhan Penduduk YoY
& 19.909 & 0.000 & 2.11 & 🟢 SIGNIFIKAN \\
Shift Index Tambang+Industri / Pertanian & Persentase Penduduk Miskin &
47.541 & 0.000 & 3.20 & 🟢 SIGNIFIKAN \\
Porsi PDRB Tambang+Industri (B+C) & Kepadatan Penduduk Kabupaten &
11.672 & 0.001 & 0.56 & 🟢 SIGNIFIKAN \\
Porsi PDRB Tambang+Industri (B+C) & Jumlah Penduduk Kabupaten (ribu
jiwa) & 11.672 & 0.001 & 1.77 & 🟢 SIGNIFIKAN \\
Porsi PDRB Tambang+Industri (B+C) & Laju Pertumbuhan Penduduk YoY &
8.796 & 0.003 & 1.65 & 🟢 SIGNIFIKAN \\
Porsi PDRB Tambang+Industri (B+C) & Persentase Penduduk Miskin & 40.577
& 0.000 & 2.93 & 🟢 SIGNIFIKAN \\
Porsi PDRB Industri Pengolahan (C) & Kepadatan Penduduk Kabupaten &
12.164 & 0.000 & 1.79 & 🟢 SIGNIFIKAN \\
Porsi PDRB Industri Pengolahan (C) & Jumlah Penduduk Kabupaten (ribu
jiwa) & 75.322 & 0.000 & 4.42 & 🟢 SIGNIFIKAN \\
Porsi PDRB Industri Pengolahan (C) & Laju Pertumbuhan Penduduk YoY &
0.322 & 0.570 & 0.90 & 🔴 TIDAK SIGNIFIKAN \\
Porsi PDRB Industri Pengolahan (C) & Persentase Penduduk Miskin & 8.061
& 0.005 & 0.62 & 🟢 SIGNIFIKAN \\
\end{longtable}
}

\begin{quote}
\textbf{Ringkasan Eksekutif:}

Sebagian skenario menunjukkan hubungan signifikan antara intensitas
ekonomi ekstraktif dan tekanan sosial-demografis. Karena unit observasi
sudah diturunkan ke kabupaten-tahun, hasil ini lebih peka terhadap
variasi lokal dibanding panel provinsi-tahun. Temuan ini memperkuat
argumen bahwa hilirisasi nikel bukan hanya fenomena ekonomi sektoral,
tetapi juga perubahan struktural yang menekan ruang hidup dan kesehatan
publik.
\end{quote}

\end{document}
